Fix an integer \(\ell\ge 1\). For a positive integer \(m\) with prime factorisation
\[m=\prod_{p} p^{v_p(m)},\]
define its \emph{powerful part}
\[Q_2(m):=\prod_{p:\,v_p(m)\ge 2} p^{v_p(m)}.\]
Equivalently, \(Q_2(m)\) is the largest powerful (squarefull) divisor of \(m\), and
\(Q_2(m)=m/\prod_{p:\,v_p(m)=1}p\).
For \(n\in\mathbb N\) set
\[P_{n,\ell}:=\prod_{i=0}^{\ell} (n+i).\]
The questions are:

(a) For every \(\varepsilon>0\) and \(\ell\ge 1\), does there exist \(N(\varepsilon,\ell)\) such that for all \(n\ge N(\varepsilon,\ell)\),
\[Q_2(P_{n,\ell})<n^{2+\varepsilon}?\]

(b) For \(\ell\ge 2\), is
\[\limsup_{n\to\infty} \frac{Q_2(P_{n,\ell})}{n^2}=+\infty?\]

(c) For \(\ell\ge 2\), is
\[\lim_{n\to\infty} \frac{Q_2(P_{n,\ell})}{n^{\ell+1}}=0?\]

Edge cases: \(n\) is intended large; \(Q_2(1)=1\). All questions are about the asymptotic regime \(n\to\infty\) with \(\ell\) fixed.
